\documentclass[../main.tex]{subfiles}

\begin{document}

    In this question I was asked to modify code to add tolerance to it.
    I then was asked to add the golden-search method to find the minimum and evaluate the log value of the error
    at each iteration.
    You can find the code in under \hyperref[subsubsec:tolerance-golden-appendix]{Appendix: Optimisation}.

    Producing these graphs:
    \begin{figure}[H]
        \begin{subfigure}
            \centering
            \includegraphics[width=0.5\textwidth]{optimisation-full-plot}
            \caption{Plot of $f(x) = e^{-(x-\pi)} + x^2$.}
            \label{fig:optimisation-full-plot}
        \end{subfigure}
        \begin{subfigure}
            \centering
            \includegraphics[width=0.5\textwidth]{optimisation-dichotomous-error}
            \caption{Logarithmic dichotomous error against iteration step.}
            \label{fig:optimisation-dichotomous-errors}
        \end{subfigure}
        \begin{subfigure}
            \centering
            \includegraphics[width=0.5\textwidth]{optimisation-golden-error}
            \caption{Logarithmic golden-section error against iteration step.}
            \label{fig:optimisation-golden-errors}
        \end{subfigure}
        \label{fig:optimisation-plots}
    \end{figure}


    The minimum occurs at $x \approx 1.839$.

    As you can see, the error decreases quicker when using the dichotomous method,
    meeting the tolerance in 17 steps opposed to 22 steps.
    This is an expected result as dichotomous error halves each
    iteration and the golden-search roughly decreases by a factor of 0.618 (1 over the golden ratio).

    \subsection*{Further analysis}\label{subsec:further-analysis}

    However, with further analysis of more factors, such as wall time and function calls,
    we get the following graphs:

    \begin{figure}[H]
        \begin{subfigure}
            \centering
            \includegraphics[width=0.5\textwidth]{optimisation-dichotomous-function-calls}
            \caption{Total dichotomous function calls for current iteration.}
            \label{fig:optimisation-dichotomous-function-calls}
        \end{subfigure}
        \begin{subfigure}
            \centering
            \includegraphics[width=0.5\textwidth]{optimisation-golden-function-calls}
            \caption{Total golden-section function calls for current iteration.}
            \label{fig:optimisation-golden-function-calls}
        \end{subfigure}
        \label{fig:optimisation-plots-further-analysis}
    \end{figure} \\

    The linear nature of these plots allows us to estimate the number of function calls per iteration
    rather easily.
    We calculate that, on average, the dichotomous method has $2.8 \approx 3$ function calls per iteration and
    the golden-search method has $0.97 \approx 1$ function call per iteration.
    This is a significant improvement in performance and is evident when measuring the wall time of each method:
    the dichotomous program runs for 6.4 milliseconds and the golden-search method runs for 0.2 milliseconds
    ($\approx 32\times$ faster).
    \\ \\
    As an experienced programmer, these results are not surprising as function calls are far more
    expensive than inline arithmetic calculations for CPUs and therefore minimising function calls should be
    prioritised over the number of iterations when designing algorithms.
    If required, function-dominated behaviour could be achieved by using \verb|time.sleep()|
    to introduce an artificial delay within the function being minimised, but I don't recommend this for large numbers
    of iterations as programs could take a long time to run.
    \\ \\
    If you're interested to know more, please read \hyperref[sec:why-are-function-calls-expensive]{Appendix: Why are function calls expensive?} for more information on this topic.

\end{document}
